\subsection{INVARIANT FIELDS OF POINT DISTRIBUTIONS ON A LIE GROUP}

DEFINITION 4. Let G be a Lie group. A field of distributions on G is called left (resp. right) invariant if it is invariant under left (resp. right) translations of G.

In other words, a field of distributions \(g \mapsto \Delta_g\) on \(G\) is left invariant if

\[
\Delta_ {g g ^ {\prime}} = \gamma (g) _ {*} \Delta_ {g ^ {\prime}} \quad \text { for } g, g ^ {\prime} \text { in } G,
\]

or again if

\[
\Delta_ {g g ^ {\prime}} = \varepsilon_ {g} * \Delta_ {g ^ {\prime}} \quad \text { for   } g, g ^ {\prime} \text {   in   G. }
\]

It is right invariant if

\[
\Delta_ {g g ^ {\prime}} = \delta (g ^ {\prime - 1}) _ {*} \Delta_ {g} \quad \text { for   } g, g ^ {\prime} \text {   in   G },
\]

or again if.

\[
\Delta_ {g g ^ {\prime}} = \Delta_ {g} * \varepsilon_ {g ^ {\prime}}, \quad \text { for   } g, g ^ {\prime} \text {   in   G. }
\]

DEFINITION 5. Let G be a Lie group and \(t \in \mathrm{U}(\mathrm{G})\) . Let \(\mathbf{L}_t\) denote the field of distributions \(g \mapsto \varepsilon_g * t\) on G and \(\mathbb{R}_t\) the field of distributions \(g \mapsto t * \varepsilon_g\) on G.

In other words, \(L_{t}\) (resp. \(R_{t}\) ) is the field of distributions defined by t and G operating on G on the right (resp. left) by means of the mapping \((g, g') \mapsto gg'\) . Let \(\Omega\) be an open subset of G and F a Hausdorff polynormed space; if \(f \in \mathcal{C}^{\omega}(\Omega, \mathrm{F})\) , then \(\mathrm{L}_{t}f = f * t^{\vee} \in \mathcal{C}^{\omega}(\Omega, \mathrm{F})\) and

\[
\mathrm{R} _ {t} f = t ^ {\vee} * f \in \mathcal {C} ^ {\omega} (\Omega , \mathrm{F})
\]

(no. 5). If G is finite-dimensional, the differential operators \(L_{t}\) and \(R_{t}\) are of class \(C^{\omega}\) (no. 5).

PROPOSITION 23. (i) The mapping \(t \mapsto \mathbf{L}_t\) (resp. \(t \mapsto \mathbf{R}_t\) ) is an isomorphism of the vector space \(\mathbf{U}(\mathbf{G})\) onto the vector field of left (resp. right) invariant distributions on \(\mathbf{G}\) .

\begin{enumerate}
\def\labelenumi{(\roman{enumi})}
\setcounter{enumi}{1}
\item
  For \(t, t'\) in U(G), \(L_{t*t'} = L_t \circ L_{t'}\) , \(R_{t*t'} = R_{t'} \circ R_t\) , \(L_t \circ R_{t'} = R_{t'} \circ L_t\) (with the abuse of notation of no. 5).
\item
  If \(\theta\) is the mapping \(g \mapsto g^{-1}\) of \(\mathbf{G}\) onto \(\mathbf{G}\) , then \(\theta(\mathbf{L}_t) = \mathbf{R}_{t^{\vee}}\) .
\item
  If \(t \in \mathrm{U}(\mathrm{G})\) and \(g \in \mathrm{G}\) , then \((\mathrm{L}_t)_g = (\mathrm{R}_{\varepsilon_g * t * \varepsilon_{g^{-1}}})_g\) .
\end{enumerate}

In G every right translation commutes with every left translation. By Proposition 21 of no. 5, \(L_{t}\) is therefore left invariant. As \((\mathbf{L}_{t})_{e}=t\) , the mapping \(t\mapsto L_{t}\) is injective. Let \(\Delta\) be a field of left invariant distributions on G; let \(t=\Delta_{e}\) ; then \(\Delta\) and \(L_{t}\) have the same value at e and are left invariant and hence \(\Delta=L_{t}\) . This proves (i) for \(L_{t}\) and the argument is similar for \(R_{t}\) . The formulae \(L_{t*t^{\prime}}=L_{t}\circ L_{t^{\prime}}\) , \(R_{t*t^{\prime}}=R_{t^{\prime}}\circ R_{t}\) follow from (21) and (18). Let \(t\in\mathrm{U}_{s}(\mathrm{G})\) , \(t^{\prime}\in\mathrm{U}_{s^{\prime}}(\mathrm{G}), f\in\mathcal{C}^{r}(\Omega,\mathrm{F})\) , where \(\Omega\) is open in G and \(s+s^{\prime}\leqslant r\) ; then

\[
\begin{array}{r l} \mathrm{L} _ {t} \mathrm{R} _ {t ^ {\prime}} f & = \mathrm{L} _ {t} (t ^ {\prime^ {\vee}} * f) = (t ^ {\prime^ {\vee}} * f) * t \\ & = t ^ {\prime^ {\vee}} * (f * t ^ {\vee}) \quad \text {(Proposition 20)} \\ & = \mathrm{R} _ {t ^ {\prime}} \mathrm{L} _ {t} f \end{array}
\]

and hence \(\mathbf{L}_t\circ \mathbf{R}_{t'} = \mathbf{R}_{t'}\circ \mathbf{L}_t\) . As \(\theta\) is an isomorphism of \(\mathbf{G}\) onto \(\mathbf{G}^{\vee}\) , \(\theta (\mathrm{L}_t)\) is a field of right invariant distributions on \(\mathbf{G}\) ; its value at \(e\) is \(\theta^{*}(t) = t^{\vee}\) hence \(\theta (\mathrm{L}_t) = \mathrm{R}_{t^{\vee}}\) . Finally,

\[
\left(\mathrm{L} _ {t}\right) _ {g} = \varepsilon_ {g} * t = \left(\varepsilon_ {g} * t * \varepsilon_ {g^{-1}}\right) * \varepsilon_ {g} = \left(\mathrm{R} _ {\varepsilon_ {g} * t * \varepsilon_ {g^{-1}}}\right) _ {g}.
\]

Remark 1. It is the action of G on itself by right translation which defines the fields of left invariant distributions.

Remark 2. Suppose that G is finite-dimensional. The mapping

\[
(t, g) \mapsto (\mathbf {R} _ {t}) _ {g} = t * \varepsilon_ {g}
\]

of \(\mathrm{U}_{s}(\mathrm{G})\times\mathrm{G}\) into \(\mathrm{T}^{(s)}(\mathrm{G})\) is an isomorphism of analytic vector bundles; for this mapping is bijective, linear on each fibre and analytic (no. 5); on the other hand, let \(\phi:\mathrm{T}^{(s)}(\mathrm{G})\to\mathrm{U}_{s}(\mathrm{G})\times\mathrm{G}\) be the inverse bijection; if \(t\in\mathrm{T}_{g}^{(s)}(\mathrm{G})\) , then \(\phi(t)=(t*\varepsilon_{g^{-1}},g)\) and hence \(\phi\) is analytic. The isomorphism \(\phi\) is called the right trivialization of \(\mathrm{T}^{(s)}(\mathrm{G})\) . Similarly, consider the mapping \((t,g)\mapsto(\mathrm{L}_{\mathrm{t}})_{g}=\varepsilon_{g}*t\) of \(\mathrm{U}_{s}(\mathrm{G})\times\mathrm{G}\) into \(\mathrm{T}^{(s)}(\mathrm{G})\) ; the inverse isomorphism is called the left trivialization of \(\mathrm{T}^{(s)}(\mathrm{G})\) . By restriction we recover the right and left trivializations of \(\mathrm{T}(\mathrm{G})\) (§2, no. 2).

